\begin{lemma}[Change lemma: both vertices in the starred stage]
Fix a finite loopless multigraph with its exact edge set, a palette of
$k$ colors, and a distinguished edge $e_0$.  Let $(T,\varphi_n)$ be a
positive-rung extended Tashkinov tree.  Let $\mathcal B$ be a hierarchy
base and let
\[
 T_n=T_{n,0}\subset T_{n,1}\subset\cdots\subset
 T_{n,q}\subset T_{n,q+1}=T
\]
be a good hierarchy under $\varphi_n$, where $q$ is its last divider
index and
\[
 A=T_{n,q}^{*}=\noderef{HierarchyStarredStage}(H,q).
\]
Assume that $T$ has maximum property, that $V(T)$ is not elementary,
that $\varphi_n$ is proper and leaves exactly $e_0$ uncolored, and that
the rung number is positive.  These data form a
\noderef{HierarchyCounterexample} $C$.

Assume the six lower-rung assertions of the technical theorem, uniformly
for every positive-rung maximum ETT on the same graph, palette, and
distinguished edge, with arbitrary certified input and generating
colorings and with rung number smaller than that of $T$: lower-rung
elementarity, terminal interchangeability, the two PE stability and
stored-pair assertions, the terminal PE component assertion, the
defective-endpoint order assertion, and the mod-coloring assertion.
Assume also the same-rung induction hypothesis that every maximum ETT
with a good hierarchy having fewer than $q$ dividers is elementary.
Finally assume the large-color estimate
\[
 2n+11\leq |\overline\varphi_n(T_n)|.
\]

Put $\ell=|T|-|A|$, the number of exact tail additions above $A$, and let
\[
 T(y_{\ell-1})=T\big|_{\min\{|A|+(\ell-1),|T|\}}
\]
be the exact ordered prefix supplied by \noderef{TreeSequencePrefix}.
Assume (6.6): $V(T(y_{\ell-1}))$ is elementary under $\varphi_n$.
Assume moreover that $C$ is lexicographically minimal among all
\noderef{HierarchyCounterexample}s on the fixed graph, palette, and
$e_0$ which have exactly $q$ dividers, allowing their certified input
colorings, generating colorings, bases, ETTs, and good hierarchies to
vary.  Explicitly, no such counterexample $C'$ satisfies
\[
 \noderef{HierarchyCounterexampleLess}
  \bigl(\noderef{HierarchyCounterexampleMeasure}(C'),
        \noderef{HierarchyCounterexampleMeasure}(C)\bigr).
\]

Let $\alpha\ne\beta$ be colors missing on $T(y_{\ell-1})$.  Write
\[
 i_c=\noderef{HierarchyFirstMissingPosition}(\varphi_n,T,c),
 \qquad v_c=T.\mathit{vertex}(i_c).
\]
If at least one of $\alpha,\beta$ is not missing on $A$, assume both
\[
 v_\alpha\prec v_\beta                                      \tag{a}
\]
and the exact edge-identity exclusion
\[
 \alpha\notin
 \left\{c:\exists t\in\operatorname{Fin}|T|,
       |A|\leq t<i_\beta,
       \ \varphi_n(T.\mathit{edge}(t))=c\right\}.             \tag{b}
\]
Thus (b) is precisely
$\alpha\notin\varphi_n\langle T(v_\beta)-A\rangle$; it is not
an endpoint-only abbreviation.


Assume in addition that both first-missing vertices lie in the last
starred stage:
\[
        v_\alpha,v_\beta\in V(A).
\]
Then the full two-color components through $v_\alpha$ and $v_\beta$
are equal in each of the following cases.
\begin{enumerate}
\item $q\geq1$, and either
$\alpha\in\overline\varphi_n(T_{n,q})$ or
$\{\alpha,\beta\}\cap D_{n,q}=\varnothing$.
\item $q=0$, and either
$\alpha\in\overline\varphi_n(T_n)$ or
$\{\alpha,\beta\}\cap D_n=\varnothing$.
\item $\alpha\in\overline\varphi_n(A)$ and the singleton
$\{\alpha\}$ is closed on $A$ under $\varphi_n$.
\end{enumerate}
In case 3 the colors are also $A$-interchangeable: at most one paper
$(\alpha,\beta)$-path component meets $A$, with an isolated vertex
component counted as a path.
\end{lemma}
\begin{proof}
Unpack $C$ through \noderef{HierarchyCounterexample}, and explicitly set
\[
 H=C.\mathit{hierarchy},\qquad
 a=v_\alpha=T.\mathit{vertex}(i_\alpha),\qquad
 b=v_\beta=T.\mathit{vertex}(i_\beta).
\]
Thus $H$ is the displayed \noderef{TashkinovHierarchy}, its good-family
witness is $C.\mathit{good}$, its maximum-property witness is
$C.\mathit{maximum}$, and $C.\mathit{notElementary}$ says that $V(T)$ is
not elementary.  For every vertex $x$, let $K_x$ denote the full
\noderef{TwoColorComponent} through the vertex-node $x$.  All components
below therefore contain the actual edge nodes of the multigraph; parallel
edges are never identified by their endpoint pairs.

We first discharge the zero-tail endpoint.  The nested-stage and terminal fields of
\noderef{TashkinovHierarchy}, together with the two cases of
\noderef{HierarchyStarredStage}, make $A$ an ordered initial segment of
$T$.  At index zero this is the `coreInsideFirst' field followed by the
stage chain, and at a positive index it is the corresponding ordinary
stage followed by that chain.  Hence
\[
                         |A|\leq |T|.                         \tag{0a}
\]

Suppose that $\ell=0$.  Since $\ell=|T|-|A|$ and (0a) gives
$|A|\leq |T|$, we have $|A|=|T|$.  Expanding
\noderef{TreeSequenceInitialSegment}, the already established inclusion
$A\subseteq T$ identifies every vertex entry and every exact edge entry;
because the lengths are equal, it identifies the ordered tree-sequences
themselves, so $A=T$.

For $\ell=0$, truncated subtraction gives $\ell-1=0$, and the length of the
displayed prefix in the statement is therefore
\[
 \min\{|A|+(\ell-1),|T|\}=\min\{|T|,|T|\}=|T|.
\]
The full-length instance of \noderef{TreeSequencePrefix} is entrywise
equal to $T$.  Thus (6.6), the assumed elementarity of this prefix, says
that $V(T)$ is elementary under $\varphi_n$, contradicting the
`notElementary' field of $C$ in
\noderef{HierarchyCounterexample}.  We have proved $\ell\ne0$, and hence
$\ell\geq1$.

It follows that $A$, every starred stage, and every ordinary hierarchy
stage through $T_{n,q}$ are ordered initial segments of
$T(y_{\ell-1})$.  Indeed, the prefix length is at least $|A|$ because
$\ell\geq1$, while \noderef{HierarchyBase},
\noderef{TashkinovHierarchy},
\noderef{HierarchyStarredStage}, and
\noderef{TreeSequenceInitialSegment} put all the named stages inside
$A$ entry by entry.  Restricting the pairwise disjoint missing-color
condition in \noderef{ColoredVertexSetElementary} therefore shows that
all these stages, including $A$, are elementary under $\varphi_n$.     \tag{1}

Because $\alpha$ and $\beta$ are missing on this elementary prefix,
the least-candidate clause of
\noderef{HierarchyFirstMissingPosition} makes $a=v_\alpha$ and
$b=v_\beta$ actual vertices of the prefix missing the indicated colors.
The location hypothesis puts both in $A$.  Hence $a$ is the unique
vertex of $A$ missing $\alpha$, and $b$ is the unique vertex of $A$
missing $\beta$.

We now establish once the paper's non-PE terminal closedness consequence.
Let $n$ be the last rung index and let $T_n$ be its terminal ladder tree.
Suppose the terminal extension kind is not PE, and let
$c\in\overline\varphi_n(T_n)$.  The \noderef{TashkinovSeries} stored in
the \noderef{ExtendedTashkinovTree} certifies its last positive step by
\noderef{TashkinovSeriesStep}; positivity excludes the initial
constructor, and the extension-kind assumption excludes PE, so the step
is SE or RE.

In the SE case, let $\varphi_{n-1}$ be the preceding coloring and let
`choice' be its maximum-defective certificate.  The `currentColoring'
field identifies $\varphi_n$ with `choice.selected'.  The
`selectedStable' field of \noderef{MaximumDefectiveChoice} says that
this selected coloring is stable on the exact tree $T_n$ relative to
$\varphi_{n-1}$ and the preceding stored set.  Expanding
\noderef{StableColoring}, its vertex clause gives
\[
 \overline{\varphi_n}(x)=\overline{\varphi_{n-1}}(x)
       \quad\text{for every }x\in V(T_n),                    \tag{2a}
\]
and hence equality of the two tree-missing sets.  The SE constructor
gives properness and exact one-uncoloredness for $\varphi_{n-1}$, while
\noderef{MaximumDefectiveChoice} gives the same two facts for
`choice.selected'.  The symmetry clause of \noderef{sc2} therefore
supplies the reverse stable relation, from $\varphi_n$ to
$\varphi_{n-1}$, with the same tree and stored set.

Choose $x\in V(T_n)$ missing $c$ under $\varphi_n$.  Equation (2a) says
that $x$ also misses $c$ under $\varphi_{n-1}$.  If an exact boundary-edge
identity $f$ of $T_n$ had color $c$ under $\varphi_n$, then $f$ would be
incident with $T_n$.  Since $c$ is tree-missing, the protected
incident-edge clause of the reverse \noderef{StableColoring} relation
would give
\[
                    \varphi_{n-1}(f)=\varphi_n(f)=c.
\]
This contradicts the `previousClosed' field of the SE constructor:
expanding \noderef{ColoredVertexSetClosed}, a boundary edge cannot have
a color missing at the inside vertex $x$.

In the RE case, the `currentColoring' field of
\noderef{TashkinovSeriesStep} identifies the current coloring literally
with the preceding coloring.  Thus the ETT generating-color equation
identifies $\varphi_n$ with the coloring for which `previousClosed'
certifies $T_n$ closed.  The definitions
\noderef{ColoredVertexSetClosed} and \noderef{BoundaryColorEdges} again
show directly that no $c$-colored boundary edge exists.  We have proved
\[
 \Theta_n\ne PE,\quad c\in\overline\varphi_n(T_n)
 \quad\Longrightarrow\quad
 T_n\text{ is }\{c\}\text{-closed under }\varphi_n.          \tag{2b}
\]
In particular, when $q=0$ and the terminal kind is not PE,
\noderef{HierarchyBase} and \noderef{HierarchyStarredStage} identify
$A=T_n$.  Since $a$ and $b$ miss $\alpha$ and $\beta$ there, respectively,
(2b) makes both colors closed on $A$.

If $a=b$, then $K_a=K_b$ immediately.  Alternative 3 additionally asks
for interchangeability.  If $q>0$, then $A=T_{n,q}$ and the last-stage
instance of \noderef{step1} applies to the missing, closed color
$\alpha$.  If $q=0$ and the terminal kind is PE, conclusion (iii) of
\noderef{interchange} applies to the exact core $A$: reflexivity gives
both clauses of \noderef{CombinedStableColoring}, while
$C.\mathit{maximum}$, $C.\mathit{positive}$ and the stated lower-rung
package give its remaining structural premises.  If $q=0$ and the
terminal kind is not PE, (2b) makes both colors closed.  Let $Q$ be any
\noderef{TwoColorPathComponent} meeting $A$.  Closedness prevents it
from crossing the boundary.  If it contains an edge-node, its path
witness, \noderef{TwoColorComponentSpanningPath}, and
\noderef{TwoColorIncidenceDegreeAtMostTwo} give two distinct graph-vertex
ends in $A$, each missing $\alpha$ or $\beta$.  Elementarity would force
both ends to equal the common unique vertex $a=b$, a contradiction.  If
$Q$ has no edge-node, its tree vertex misses both colors by
\noderef{MissingColors}; elementarity identifies that vertex with
$a=b$, so $Q=K_a$.  Hence $K_a$ is the unique path component meeting
$A$, which is exactly \noderef{TreeInterchangeable}.  Thus the
coincident-vertex branch is complete without any appeal to a truncated
lower-rung ETT.

We may henceforth assume $a\ne b$.  Elementarity says that $\beta$,
which is missing at $b$, is not missing at $a$, and symmetrically
$\alpha$ is not missing at $b$.  By the complementary definitions
\noderef{MissingColors} and \noderef{PresentColors}, each vertex is
incident with an edge of the other color.  Properness makes it a
degree-one end in the two-color incidence graph.  Hence $K_a$ and $K_b$
are nontrivial path components in the sense of
\noderef{TwoColorNontrivialPathComponent}.  This verifies the explicit
\noderef{TwoColorPathComponent} premise each time
\noderef{TreeInterchangeable} is used below.                             \tag{2}

We shall repeatedly use the following endpoint observation, always with
its path hypothesis explicit.  Let $B$ be an elementary hierarchy stage
which contains $a$ and $b$, and let $Q$ be a
\noderef{TwoColorPathComponent} meeting $B$.  If both $\alpha$ and
$\beta$ are closed on $B$, then $Q$ cannot cross the boundary of $B$:
the first incidence step from a graph vertex in $B$ to one outside it
would use an exact boundary-edge identity colored $\alpha$ or $\beta$,
contrary to \noderef{TreeColorSetClosed} and
\noderef{BoundaryColorEdges}.

Suppose first that $Q$ contains an edge-node.  The leaf witness in
\noderef{TwoColorPathComponent}, together with the spanning path from
\noderef{TwoColorComponentSpanningPath} and the degree bound
\noderef{TwoColorIncidenceDegreeAtMostTwo}, makes $Q$ a nontrivial path
whose two spanning-path ends have incidence degree one.  They are
graph-vertex nodes: a colored edge-node has the two distinct endpoint
neighbors because the graph is loopless.  Properness permits at most one
incident edge of either color at a graph vertex, so each degree-one end
misses at least one of $\alpha,\beta$.  Both ends lie in $B$.  They cannot
both miss the same color, because $B$ is elementary.  Hence they are the
unique $\alpha$-missing and $\beta$-missing vertices $a$ and $b$, and
$Q=K_a=K_b$.

Suppose instead that $Q$ has no edge-node.  Choose $x\in V(B)$ whose
vertex-node lies in $Q$.  By \noderef{TwoColorIncidenceGraph}, any exact
edge incident with $x$ and colored $\alpha$ or $\beta$ would be adjacent
to $x$ and hence, by \noderef{TwoColorComponent}, its edge-node would
belong to $Q$.  Thus neither color is present at $x$, so
\noderef{MissingColors} and \noderef{PresentColors} show that both are
missing there.  Elementarity forces $x=a$ and $x=b$, and again
$Q=K_a=K_b$.  Applying this conclusion to two arbitrary path components
meeting $B$ proves
\[
 \begin{gathered}
 B\text{ elementary},\quad a,b\in V(B),\quad
 B\text{ is }\{\alpha\}\text{- and }\{\beta\}\text{-closed}\\
 \Longrightarrow\quad K_a=K_b\text{ and }\alpha,\beta
     \text{ are }B\text{-interchangeable}.
 \end{gathered}                                                         \tag{3}
\]
This observation makes no endpoint claim for a cycle: a cycle is not a
\noderef{TwoColorPathComponent} and does not enter the definition of
\noderef{TreeInterchangeable}.

\textbf{Location 1: $a,b\in V(A)$.}
By (1), $A$ is elementary; hence $a$ and $b$ are respectively the
unique vertices of $A$ missing $\alpha$ and $\beta$.  If both singleton
colors are closed on $A$, apply (3) with $B=A$.  It gives $K_a=K_b$ and
also proves that this is the unique paper path component meeting $A$,
including among isolated-vertex paths.                                   \tag{6}

Suppose precisely one of the two colors is closed on $A$.  If $q>0$,
then $A=T_{n,q}$ and the last-divider instance of
\noderef{step1}, with the closed color placed in its first color slot,
applies.  Its lower-rung premise is exactly the six-assertion hypothesis
of this lemma, its fewer-divider premise is our same-rung induction
hypothesis, and its numerical premise is the displayed $2n+11$ estimate.
It makes the two colors $A$-interchangeable, placing whichever color is
closed in the first slot.  If $q=0$ and the terminal
kind is PE, use the core conclusion (iii) of
\noderef{interchange}; reflexivity of the two clauses in
\noderef{CombinedStableColoring} supplies its stability premise, and
our six lower-rung package supplies its induction premise.  If $q=0$
and the terminal kind is not PE, the already proved consequence (2b)
makes both colors closed on $A=T_n$.  Thus the assumed ``precisely one''
subcase cannot occur.  The definition \noderef{TwoColorIncidenceGraph}
is symmetric in $\alpha,\beta$, so interchanging their slots does not
change any component or the meaning of interchangeability.  In the two
possible cases, (2) says explicitly
that $K_a$ and $K_b$ are path components meeting $A$; applying the
resulting interchangeability through the definition
\noderef{TreeInterchangeable} gives $K_a=K_b$ and uniqueness of the path
component meeting $A$.                                                 \tag{7}

It remains in this location to consider the case in which neither color
is closed on $A$.  Before choosing a first missing stage, we justify the
paper's reduction $a\prec b$.  The strict position order is the one in
\noderef{TreeVertexOrder}; distinct displayed vertices have distinct
indices by \noderef{TreeSequence}.  If $q>0$, then
$A=T_{n,q}$ by \noderef{HierarchyStarredStage}, so both named colors are
missing on $T_{n,q}$.  The local hypotheses and the desired component
equality are therefore symmetric in the two colors, and in the order
$b\prec a$ we interchange the names $(\alpha,a)$ and $(\beta,b)$.  If
$q=0$ and $\{\alpha,\beta\}\cap D_n=\varnothing$, the same relabeling is
legitimate because pair avoidance is symmetric.

There remains the asymmetric $q=0$ alternative in which pair avoidance
fails and $\alpha$ is specifically missing on $T_n$.  The definition
\noderef{TashkinovStoredColors} first gives $\alpha\notin D_n$; failure of
pair avoidance consequently gives $\beta\in D_n$, and the same
definition gives $\beta\notin\overline\varphi_n(T_n)$.  Hence the first
$\alpha$-missing vertex $a$ lies in $T_n$, whereas the first
$\beta$-missing vertex $b$, which lies in $A$, is outside $T_n$.  As
the non-PE clause of \noderef{HierarchyBase} would identify $A$ with
$T_n$, this already forces the terminal kind to be PE.  The PE clause of
the same definition and the left-segment clause of
\noderef{OrderedTreeUnion} make $T_n$ an exact ordered initial segment
of $A$; \noderef{TashkinovHierarchy} makes $A$ an ordered initial segment
of $T$.  Thus the displayed index of $a$ is strictly smaller than that
of $b$, namely $a\prec b$.  This proves the order reduction in every
asymmetric alternative without exchanging the distinguished hypothesis.

Now let $r$ be the least index for which the first $\beta$-missing
vertex $b$ belongs to $T_{n,r}^{*}$.  Such an index exists because
$b\in A=T_{n,q}^{*}$.  The exact initial-segment relations in
\noderef{TashkinovHierarchy} and
\noderef{TreeSequenceInitialSegment}, together with $a\prec b$, put $a$
in this same starred stage; thus both $a$ and $b$ lie there, as required
before the descent.  If $r>0$, the $\beta$ witness is new above the
preceding starred stage; (1) and
\noderef{GoodHierarchyNewMissingClosed} make $\beta$ closed on
$T_{n,r}=T_{n,r}^{*}$.  If $\alpha$ is also closed there, apply the
path-qualified observation (3) to this elementary stage.  Otherwise the
stage-$r$ instance of \noderef{step1} gives interchangeability.  In the
latter case, (2) supplies the path-component hypotheses for $K_a,K_b$,
and \noderef{TreeInterchangeable} gives $K_a=K_b$.

Suppose $r=0$.  Then $a,b$ both lie in the base core
$T_{n,0}^{*}$.  If the last extension is non-PE,
\noderef{HierarchyBase} identifies this core with $T_n$.  Both colors
are tree-missing there, so the already established SE/RE consequence
(2b) makes both closed.  Observation (3), applied to this elementary
base stage, gives $K_a=K_b$.

If the last extension is PE, its base clause gives the exact union
\[
                         T_{n,0}^{*}=T_n\vee R_n.
\]
The vertex-set clause of \noderef{OrderedTreeUnion} places each of $a,b$
in $T_n$ or $R_n$.  If both lie in $T_n$, conclusion (ii) of
\noderef{interchange}, applied to the color $\alpha$ missing at $a$,
makes their path components equal.  If both lie in $R_n$, conclusion (i)
does the same.  If $a\in T_n$ and $b\in R_n$, conclusion (iv) applies to
$\alpha,\beta$; if $a\in R_n$ and $b\in T_n$, apply (iv) after exchanging
the two colors.  In each use, reflexivity gives
\noderef{CombinedStableColoring}; the other premises are
$C.\mathit{maximum}$, $C.\mathit{positive}$, properness, exact
one-uncoloredness, and the stated lower-rung package.  The two components
are path components by (2), so the asserted interchangeability really
does imply $K_a=K_b$.  This is the paper's descent to the first starred
stage; it uses the PE base interface and the exact union rather than
pretending that the union itself is a ladder ETT.                       \tag{8}

Equations (6)--(8) prove the required component equality in Location 1
under each of alternatives 1--3 in the statement.  They also prove the
additional $A$-interchangeability whenever alternative 3 holds.
\end{proof}
